\subsection{Procedimentos de Execução e Coleta de Dados}

Para assegurar a validade interna e o isolamento dos fatores, o experimento foi executado em um fluxo estritamente controlado. O pipeline de execução, detalhado na Figura \ref{fig:rag_experimental_pipeline}, ilustra o percurso sistemático de cada consulta do \textit{Golden Set} através dos módulos de recuperação híbrida, fusão de \textit{rankings} e inferência final.

\begin{figure}[ht]
  \centering
  \includesvg[width=1\columnwidth]{figuras/rag}
  \caption{Pipeline de execução experimental do sistema RAG.}
  \label{fig:rag_experimental_pipeline}
\end{figure}

\subsubsection{Controle de Variáveis}

A fim de neutralizar variáveis de confusão, a temperatura dos modelos geradores foi fixada em $T = 0.1$ e utilizou-se um único \textit{System Prompt} padronizado (Apêndice B) com restrições negativas contra alucinações.

\subsubsection{Fluxo de Inferência}

Conforme o esquema detalhado na Figura \ref{fig:rag_experimental_pipeline}, a execução experimental ocorreu em duas fases sequenciais e interdependentes:

\paragraph{Fase 1: Recuperação (\textit{Retrieval})}
Esta fase inicia-se no Banco de Perguntas, onde cada consulta é direcionada ao módulo de Busca. O sistema opera em três frentes paralelas: a similaridade de cosseno sobre o \textit{ChromaDB} (Busca Vetorial), o algoritmo BM25 (Busca Lexical) ou a integração de ambos via \textit{Reciprocal Rank Fusion} (RRF). O resultado desta etapa é a consolidação de um Banco de Contexto contendo os documentos recuperados com maior relevância estatística.

\paragraph{Fase 2: Geração (\textit{Generation})}
O módulo de geração consome os dados do Banco de Contexto sob um critério de seleção determinística (\textit{Top-k}), simulando janelas de contexto com $\{5, 10, 15, 20\}$ fragmentos.  Os textos selecionados são injetados nos modelos \texttt{gpt-4o} ou \texttt{gpt-4o-mini}, conforme o tratamento experimental. A resposta gerada é então persistida em um Banco de Respostas tabular (CSV) para análise posterior.

\subsubsection{Protocolo de Avaliação (\textit{LLM-as-a-Judge})}

Com o \textit{dataset} de inferências concluído, procedeu-se à mensuração utilizando o \textit{framework Ragas}. Instanciou-se um juiz sintético (\texttt{gpt-4o}, $T=0$) para garantir determinismo absoluto na avaliação.

O avaliador consumiu exclusivamente o contexto efetivamente injetado no \textit{prompt} após a etapa de fatiamento do \textit{Top-k}, assegurando que as métricas de fidelidade e precisão refletissem com exatidão a relação entre o insumo recuperado e o produto final gerado para cada uma das 3.840 unidades experimentais.
